\documentclass[11pt,aspectratio=169]{beamer}
\usepackage[T1]{fontenc}
\usepackage{lmodern}
\usepackage{amsmath,amssymb}
\usepackage{booktabs,tabularx,array}
\setbeamertemplate{navigation symbols}{}
\setbeamertemplate{itemize items}[circle]
\setbeamertemplate{footline}{\hfill\scriptsize PREVIEW\quad\insertframenumber/\inserttotalframenumber\hspace{1em}\vspace{0.5em}}
\setbeamercolor{frametitle}{fg=blue!55!black}
\linespread{1.2}
\newcolumntype{Y}{>{\raggedright\arraybackslash}X}
\newcommand{\src}[1]{\vfill{\tiny\color{gray}Source: #1}}
\begin{document}

\begin{frame}{Family space and household choices}
\small
Physical rooms $h$ supply housing space; $m$ is the number of children currently at home.
When $m>0$, the saved working economy requires $h>h_P=2.594$ rooms.
\par\medskip
\begin{center}
\begin{tabular}{lll}
\toprule
Household & Two-room space admissibility & Smallest admissible owner home\\
\midrule
No children at home & Admissible & Two rooms\\
Children at home & Inadmissible & Four rooms\\
\bottomrule
\end{tabular}
\end{center}
\par\medskip
Households choose housing, tenure, consumption and saving within the same budget; rental housing also supplies family space.
\src{Working chain 13; purchase-access census, October 4, 2026.}
\end{frame}

\begin{frame}{Housing costs and fertility}
\small
In a saved fixed-price lifecycle comparison, imposed house prices and rents rise by 10\%.
\par\medskip
\begin{center}
\begin{tabular}{lr}
\toprule
Model outcome & Change\\
\midrule
Explicit births & $-6.434\%$\\
Ownership, ages 18--29 & $-3.818$ percentage points\\
Rooms occupied, ages 18--29 & $-8.605\%$\\
\bottomrule
\end{tabular}
\end{center}
\par\medskip
Higher housing costs substantially reduce births in this comparison. The price change also affects collateral and incumbent housing wealth; it does not isolate the physical room floor.
\src{Matched price diagnostic, October 4, 2026; imposed prices, not a dated transition.}
\end{frame}

\begin{frame}{Credit and the value of parenthood}
\small
The financed share of a house price rises from 80\% to 95\% at imposed prices.
\par\medskip
\begin{center}
\begin{tabular}{lr}
\toprule
Model outcome & Change\\
\midrule
Explicit births & $-0.250\%$\\
Ownership, ages 18--29 & $+10.112$ percentage points\\
\bottomrule
\end{tabular}
\end{center}
\par\medskip
The birth decision depends on the value after a successful birth relative to waiting. In an audited age-22 renter slice covering 21.5\% of first-birth responsiveness, borrowing relief raises the value of waiting more on average; signs differ by income.
\src{Matched credit diagnostic and joint-budget age-22 renter audit, October 4, 2026.}
\end{frame}

\begin{frame}{Borrowing without collateral constraints}
\small
\textbf{Proposed comparison --- no numerical results yet.}
\par\medskip
Remove owner origination and collateral restrictions and renter borrowing limits, while retaining explicit solvency and terminal-debt conditions.
\par\medskip
Hold income risk, the parent space floor and rental-size restrictions fixed. Compare births and ownership at common prices and household states first.
\par\medskip
Setting the financed share to one removes a nominal down payment; it is not an unconstrained-borrowing economy.
\src{Proposed design based on the October 4 credit-mechanism memo.}
\end{frame}

\begin{frame}{Fertility decline with a fixed housing stock}
\small
\textbf{Proposed comparison --- no numerical results yet.}
\par\medskip
Begin both transitions from the same 2007 state and retain the same fertility-preference shock. Compare elastic housing supply with a constant physical stock $\bar H=H_{2007}$.
\par\medskip
The outcomes are fertility and unit housing costs along the dated transition. Fixing a supply coefficient $H_0$ would not hold the physical stock fixed.
\par\medskip
A 2023 initial state would be a separate prospective exercise.
\src{Proposed transition design; no transition or new stationary result is asserted.}
\end{frame}
\end{document}
